\section{Оптимальный алгоритм}\label{sec:fr_optim}

Данный алгоритм характеризуется тем, что в серии из $K$ изменений периода ${T_{z}}$ инжектируемого сигнала (${T^{(k)}_{z}}$, ${k \in \left[1;K \right]}$):
\begin{itemize}
    \item Длительность ${T_{h}}$ измерительного окна кратен периоду инжектируемого сигнала (${T^{(k)}_{h} = L^{(k)}T^{(k)}_{z}}$, ${T^{(k)}_{h}}$ = var).
    \item Период дискретизации ${T_{s}}$ кратен периоду ${T_{d}}$ выполнения операций в системе (${T^{(k)}_{s} = M^{(k)} \cdot T_d}$); коэффициент децимации ${M^{(k)} = var}$, ${M \in \mathbb{N}_{> 0}}$.
    \item Количество точек ${N \in \mathbb{N}_{> 0}}$ дискретной последовательности, по которым вычисляются коэффициенты разложения в выражении \eqref{eqn:DFT}, соответствует конкретному значению ${T^{(k)}_{z}}$ (${N^{(k)} = var}$).
    \begin{equation*}
        N^{(k)} = \cfrac{T^{(k)}_{h}} {T^{(k)}_{s}}.
    \end{equation*}
\end{itemize}

Соотношение между периодом инжектируемого сигнала и периодом дискретизации
\begin{equation*}
        R^{(k)} = \cfrac{T^{(k)}_{z}} {T^{(k)}_{s}},
\end{equation*}
где ${R^{(k)} \in \mathbb{R}}$

Длительность измерительного окна
\begin{equation*}
    T^{(k)}_{h} = L^{(k)} \cdot T^{(k)}_{z} = L^{(k)} \cdot R^{(k)} \cdot T^{(k)}_{s} = \underbrace{L^{(k)} \cdot R^{(k)}}_{N^{(k)}} \cdot \underbrace{M^{(k)} \cdot T_d}_{T^{(k)}_{s}}.
\end{equation*}

Таким образом, для каждой $k$-й частоты ${f_{z}}$ задача сводится к поиску набора целых чисел $\mathbf{p}^{(k)} = \left(N^{(k)}, M^{(k)}, L^{(k)}\right) \in \mathbb{N}_{> 0}^3$, при которых обеспечивается частота ${f^{\*}_{z}}$ с заданной степенью приближения к ${f_{z}}$.


При ограничениях
\begin{equation*}
    N_\text{min} \leq N^{(k)} \leq N_\text{max},
\end{equation*}


Если несколько наборов параметров удовлетворяют заданной точности, то из них выбирается с наименьшим $N$.

Если ни один набор параметров не удовлетворят заданной точности, то выбирается набор с наименьшим отклонением ${f^{\*}_{z}}$ от заданной ${f_{z}}$.

Последовательность в алгоритме:
\begin{enumerate}
    \item Вычисляется коэффициент
    \item Перебираются значения $N$ из диапазона ${N_\text{min}}$, ${N_\text{max}}$.
\end{enumerate}